% !TeX program = xelatex
% UTF-8; compile twice with XeLaTeX. No external images or data files needed.
\documentclass[a4paper,landscape,10pt]{article}
\usepackage[margin=14mm,headheight=14pt,headsep=5mm,footskip=7mm]{geometry}
\usepackage{fontspec}
\setmainfont{DejaVu Sans}
\newfontfamily\ipafont{DejaVu Sans}
\usepackage[table]{xcolor}
\usepackage{array,longtable,booktabs,fancyhdr}
\definecolor{ink}{HTML}{18384D}
\definecolor{band}{HTML}{E7EFF4}
\pagestyle{fancy}\fancyhf{}
\fancyhead[L]{\small\color{ink}Parlez Vous QSO?! · Español / Spanish}
\fancyhead[R]{\small Englische Phrase · Traducción española · IPA}
\fancyfoot[L]{\scriptsize Vorlage: Spanish v 1.1 · Credits: Vicente -- EA5OO; Colin Bass (ohne Rufzeichen)}
\fancyfoot[R]{\small\thepage}
\setlength{\parindent}{0pt}
\setlength{\tabcolsep}{3mm}
\renewcommand{\arraystretch}{1.25}
\newcommand{\ipa}[1]{{\ipafont[#1]}}
\begin{document}
{\LARGE\bfseries\color{ink}Parlez Vous QSO?!}\hfill{\large Español / Spanish}

\vspace{2mm}
{\small QSO-Sprachführer mit spanischer IPA-Lautschrift. Die Lautschrift bezieht sich auf die spanische Spalte.}

\vspace{2mm}
{\footnotesize \textbf{Leseschlüssel:} \ipa{ˈ} Hauptbetonung, \ipa{ˌ} Nebenbetonung;
\ipa{θ} wie englisches „th“ in „think“, \ipa{x} wie deutsches „ch“ in „Bach“, \ipa{ɲ} wie spanisches „ñ“.
\ipa{ɾ} ist ein kurzer, \ipa{r} ein gerollter r-Laut. Grundlage: Spanien mit Unterscheidung von s und z/c.}

\vspace{1mm}
\fontsize{9}{11.8}\selectfont
\begin{longtable}{@{}>{\raggedright\arraybackslash}p{13mm}>{\raggedright\arraybackslash}p{73mm}>{\raggedright\arraybackslash}p{75mm}>{\raggedright\arraybackslash}p{90mm}@{}}
\toprule
\textbf{Ref.} & \textbf{English phrase} & \textbf{Español} & \textbf{Lautschrift (IPA)}\\
\midrule
\endfirsthead
\toprule
\textbf{Ref.} & \textbf{English phrase} & \textbf{Español} & \textbf{Lautschrift (IPA)}\\
\midrule
\endhead
\midrule
\multicolumn{4}{r}{\scriptsize Fortsetzung auf der nächsten Seite}\\
\endfoot
\bottomrule
\endlastfoot
\rowcolor{band} 1.0 & \textbf{Opening Greetings} & \multicolumn{2}{l}{\textbf{Saludos de Comienzo}} \\*
1.1 & I am very pleased to meet you. & Estoy muy contento de conocerte & \ipa{esˈtoj muj konˈtento de konoˈθeɾte} \\
1.2 & Thank you for your call & Gracias por tu llamada. & \ipa{ˈɡɾaθjas poɾ tu ʝaˈmada} \\
1.3 & Good morning & Buenos días. & \ipa{ˈbwenos ˈdi.as} \\
1.31 & Good day & Buenos días. & \ipa{ˈbwenos ˈdi.as} \\
1.32 & Good afternoon & Buenas tardes (de 12 a 18 horas). & \ipa{ˈbwenas ˈtaɾdes (de ˈdoθe a djeˈθotʃo ˈoɾas)} \\
1.33 & Good evening & Buenas tardes (después de las 18 horas). & \ipa{ˈbwenas ˈtaɾdes (desˈpwes de las djeˈθotʃo ˈoɾas)} \\
1.4 & I know only a little Spanish but I enjoy QSOs in Spanish. & Conozco solo un poco de español pero me gustan los QSOs en español. & \ipa{koˈnoθko ˈsolo un ˈpoko de espaˈɲol ˈpeɾo me ˈɡustan los ku ˈese os en espaˈɲol} \\
\rowcolor{band} 2.0 & \textbf{Operator} & \multicolumn{2}{l}{\textbf{Operador}} \\*
2.1 & My name is Peter. & Me llamo Peter. & \ipa{me ˈʝamo ˈpiːtə} \\
2.2 & I have been licensed since 1960. & Tengo licencia desde el año 1960. & \ipa{ˈtenɡo liˈθenθja ˈdesde el ˈaɲo mil nobeˈθjentos seˈsenta} \\
2.3 & I have been licensed for 43 years. & Tengo licencia desde hace 43 años. & \ipa{ˈtenɡo liˈθenθja ˈdesde ˈaθe kwaˌɾenta i ˈtɾes ˈaɲos} \\
2.4 & I live in Sittingbourne. & Vivo en Sittingbourne. & \ipa{ˈbibo en ˈsɪtɪŋbɔːn} \\
2.5 & Sittingbourne is a village. & Sittingbourne es una aldea. & \ipa{ˈsɪtɪŋbɔːn es ˈuna alˈde.a} \\
2.51 & Sittingbourne is a town. & Sittingbourne es un pueblo. & \ipa{ˈsɪtɪŋbɔːn es un ˈpweblo} \\
2.52 & Sittingbourne is a city. & Sittingbourne es una ciudad. & \ipa{ˈsɪtɪŋbɔːn es ˈuna θjuˈdad} \\
2.6 & Sittingbourne is between London and Dover. & Sittingbourne está entre Londres y Dover. & \ipa{ˈsɪtɪŋbɔːn esˈta ˈentɾe ˈlondɾes i ˈdəʊvə} \\
2.7 & I live in the South of England. & Vivo en el Sur de Inglaterra. & \ipa{ˈbibo en el suɾ de inɡlaˈtera} \\
2.71 & I live in the North of England. & Vivo en el Norte de Inglaterra. & \ipa{ˈbibo en el ˈnoɾte de inɡlaˈtera} \\
2.72 & I live in the East of England. & Vivo en el Este de Inglaterra. & \ipa{ˈbibo en el ˈeste de inɡlaˈtera} \\
2.73 & I live in West of England. & Vivo en el Oeste de Inglaterra. & \ipa{ˈbibo en el oˈeste de inɡlaˈtera} \\
2.74 & I live in South East England. & Vivo en el Sudeste de Inglaterra. & \ipa{ˈbibo en el suˈdeste de inɡlaˈtera} \\
2.75 & I live in South West England. & Vivo en el Sudoeste de Inglaterra. & \ipa{ˈbibo en el sudoˈeste de inɡlaˈtera} \\
2.76 & I live in North East England. & Vivo en el Noreste de Inglaterra. & \ipa{ˈbibo en el noˈɾeste de inɡlaˈtera} \\
2.77 & I live in North West England. & Vivo en el Noroeste de Inglaterra. & \ipa{ˈbibo en el noɾoˈeste de inɡlaˈtera} \\
2.78 & I live in central England. & Vivo en el Centro de Inglaterra. & \ipa{ˈbibo en el ˈθentɾo de inɡlaˈtera} \\
2.8 & I live in the country. & Vivo en el campo. & \ipa{ˈbibo en el ˈkampo} \\
2.81 & I live by the sea. & Vivo cerca del mar. & \ipa{ˈbibo ˈθeɾka del maɾ} \\
2.82 & I live near mountains. & Vivo cerca de las montañas. & \ipa{ˈbibo ˈθeɾka de las monˈtaɲas} \\
2.9 & Sittingbourne is 50Km from London. & Sittingbourne está a 50 Km. de Londres. & \ipa{ˈsɪtɪŋbɔːn esˈta a θinˈkwenta kiˈlometɾos de ˈlondɾes} \\
2.11 & I am retired. & Estoy jubilado. & \ipa{esˈtoj xubiˈlado} \\
2.12 & My home call sign is M0ECS & Mi indicativo es M0ECS. & \ipa{mi indikaˈtibo es ˈeme ˈθeɾo e θe ˈese} \\
\rowcolor{band} 3.0 & \textbf{Station} & \multicolumn{2}{l}{\textbf{Estación}} \\*
3.1 & My radio is an Icom 706. & Mi radio es un Icom 706. & \ipa{mi ˈradjo es un ˈaɪkɒm seteˈθjentos ˈsejs} \\
3.2 & My station is all homebrew (homemade). & Mi estación de radio es toda autoconstruida. & \ipa{mi estaˈθjon de ˈradjo es ˈtoda autokonstɾuˈida} \\
3.3 & I am running 5 Watts. & Estoy saliendo con 5 Watts. & \ipa{esˈtoj saˈljendo kon ˈθinko wats} \\
3.4 & My antenna is a long wire. & Mi antena es de hilo largo. & \ipa{mi anˈtena es de ˈilo ˈlaɾɡo} \\
3.41 & My antenna is a dipole. & Mi antena es de dipolo. & \ipa{mi anˈtena es de diˈpolo} \\
3.42 & My antenna is a G5RV. & Mi antena es de G5RV. & \ipa{mi anˈtena es de xe ˈθinko ˈere ˈube} \\
3.43 & My antenna is a doublet. & Mi antena es de doblete. & \ipa{mi anˈtena es de doˈblete} \\
3.44 & My antenna is a delta loop. & Mi antena es de delta loop. & \ipa{mi anˈtena es de ˈdelta luːp} \\
3.45 & My antenna is a vertical. & Mi antena es de vertical. & \ipa{mi anˈtena es de beɾtiˈkal} \\
3.46 & My antenna is a log periodic. & Mi antena es de logarítmica periódica. & \ipa{mi anˈtena es de loɡaˈɾitmika peˈɾjodika} \\
3.47 & My antenna is a Yagi beam. & Mi antena es de directiva Yagi. & \ipa{mi anˈtena es de diɾekˈtiba ˈʝaɡi} \\
3.48 & My antenna is a 5 element beam. & Mi antena es de directiva de 5 elementos. & \ipa{mi anˈtena es de diɾekˈtiba de ˈθinko eleˈmentos} \\
\rowcolor{band} 4.0 & \textbf{Conditions} & \multicolumn{2}{l}{\textbf{Condiciones}} \\*
4.1 & Thank you for my 599 report. & Gracias por mi informe, 599. & \ipa{ˈɡɾaθjas poɾ mi inˈfoɾme ˈθinko ˈnwebe ˈnwebe} \\
4.2 & The band conditions are good. & Las condiciones de propagación en la banda son buenas. & \ipa{las kondiˈθjones de pɾopaɡaˈθjon en la ˈbanda son ˈbwenas} \\
4.3 & The band conditions are bad. & Las condiciones de propagación en la banda son malas. & \ipa{las kondiˈθjones de pɾopaɡaˈθjon en la ˈbanda son ˈmalas} \\
4.4 & There is strong interference here. & Aquí tengo una fuerte interferencia. & \ipa{aˈki ˈtenɡo ˈuna ˈfweɾte inteɾfeˈɾenθja} \\
4.5 & Can you change to the 80m band? & ¿Puedes cambiar a la banda de 80 m.? & \ipa{ˈpwedes kamˈbjaɾ a la ˈbanda de oˈtʃenta ˈmetɾos} \\
4.6 & Please change frequency to 3.558 MHz. & Por favor cambia a la frecuencia de 3.558 MHz. & \ipa{poɾ faˈboɾ ˈkambja a la fɾeˈkwenθja de tɾes ˈkoma ˈθinko ˈθinko ˈotʃo ˌmeɡaˈeɾθjos} \\
\rowcolor{band} 5.0 & \textbf{General Queries and Comments} & \multicolumn{2}{l}{\textbf{Preguntas y comentarios generales}} \\*
5.1 & What does <\#\#\#\#\#> mean please? & ¿Qué significa <\#\#\#\#\#> ? & \ipa{ke siɡniˈfika …} \\
5.2 & Please wait a moment, I need to check my Spanish dictionary. & Espera un momento por favor, necesito consultar mi diccionario de español. & \ipa{esˈpeɾa un moˈmento poɾ faˈboɾ neθeˈsito konsulˈtaɾ mi dikθjoˈnaɾjo de espaˈɲol} \\
5.3 & How do you say, <\#\#\#\#\#> in Spanish ? & ¿Cómo dirías <\#\#\#\#\#> en español? & \ipa{ˈkomo diˈɾi.as … en espaˈɲol} \\
5.5 & My favourite modes of operation are CW / RTTY / PSK / MT63 / Olivia / Hellschreiber / Digital Modes / and SSB & Mis modos de operación favoritos son CW / RTTY / PSK / MT63 / Olivia / Hellschreiber / Modos Digitales / y SSB. & \ipa{mis ˈmodos de opeɾaˈθjon faboˈɾitos son θe ˈdoble ˈube / ˈere te te i ˈɡɾjeɡa / pe ˈese ka / ˈeme te seˌsenta i ˈtɾes / oˈlibja / ˈhɛlˌʃʁaɪbɐ / ˈmodos dixiˈtales / i ˈese ˈese be} \\
5.6 & My best DX is Australia. & Mi mejor DX es Australia. & \ipa{mi meˈxoɾ de ˈekis es awsˈtɾalja} \\
5.7 & I enjoy working contests. & Me gusta participar en los concursos. & \ipa{me ˈɡusta paɾtiθiˈpaɾ en los konˈkuɾsos} \\
5.8 & I like working QRP. & Me gusta transmitir en QRP. & \ipa{me ˈɡusta tɾansmiˈtiɾ en ku ˈere pe} \\
\rowcolor{band} 6.0 & \textbf{The Weather} & \multicolumn{2}{l}{\textbf{El Tiempo}} \\*
6.1 & It is sunny. & Hay mucho sol. & \ipa{aj ˈmutʃo sol} \\
6.2 & It is windy. & Es ventoso. & \ipa{es benˈtoso} \\
6.3 & It is raining. & Está lloviendo. & \ipa{esˈta ʝoˈbjendo} \\
6.4 & It is snowing. & Está nevando. & \ipa{esˈta neˈbando} \\
6.5 & It is very hot. & Hace mucho calor & \ipa{ˈaθe ˈmutʃo kaˈloɾ} \\
6.6 & It is very cold. & Hace mucho frío. & \ipa{ˈaθe ˈmutʃo ˈfɾi.o} \\
6.7 & The temperature is 8C. & La temperatura es 8º C. & \ipa{la tempeɾaˈtuɾa es ˈotʃo ˈɡɾados ˈθelsjus} \\
6.8 & It is overcast. & Está nublado. & \ipa{esˈta nuˈblado} \\
6.9 & It is foggy. & Hay niebla. & \ipa{aj ˈnjebla} \\
6.11 & It has been raining all day. & Llovió durante todo el día & \ipa{ʝoˈbjo duˈɾante ˈtodo el ˈdi.a} \\
6.12 & There is 20cm of snow here. & Aquí hay 20 cm. de nieve. & \ipa{aˈki aj ˈbejnte θenˈtimetɾos de ˈnjebe} \\
\rowcolor{band} 8.0 & \textbf{QSL} & \multicolumn{2}{l}{\textbf{QSL}} \\*
8.1 & My QSL card is sure via the bureau & Te envío mi tarjeta QSL con seguridad vía bureau (asociación). & \ipa{te enˈbi.o mi taɾˈxeta ku ˈese ˈele kon seɡuɾiˈdad ˈbi.a byˈʁo (asoθjaˈθjon)} \\
8.2 & QSL via eQSL & QSL vía eQSL. & \ipa{ku ˈese ˈele ˈbi.a i ku ˈese ˈele} \\
8.3 & No eQSL please & Por favor no envíe eQSL. & \ipa{poɾ faˈboɾ no enˈbi.e i ku ˈese ˈele} \\
8.4 & I will send you my card direct. & Te enviaré mi tarjeta directamente. & \ipa{te embjaˈɾe mi taɾˈxeta diɾektaˈmente} \\
8.5 & What is your address? & ¿Cuál es tu dirección? & \ipa{kwal es tu diɾekˈθjon} \\
8.6 & Please send your QSL card direct. & Por favor, envíame directamente tu tarjeta QSL. & \ipa{poɾ faˈboɾ enˈbi.ame diɾektaˈmente tu taɾˈxeta ku ˈese ˈele} \\
8.7 & My address is... & Mi dirección es… & \ipa{mi diɾekˈθjon es …} \\
8.8 & My address is shown on QSL.com / EHam.net & Mi dirección está en QSL.com / EHam.net. & \ipa{mi diɾekˈθjon esˈta en ku ˈese ˈele ˈpunto kom / i hæm ˈpunto net} \\
8.9 & Sorry I do not send QSL cards. & Lo siento, no envío tarjetas QSL. & \ipa{lo ˈsjento no enˈbi.o taɾˈxetas ku ˈese ˈele} \\
8.11 & Please send QSL card via G1JQH. & Por favor, envía tu tarjeta QSL vía G1JQH. & \ipa{poɾ faˈboɾ enˈbi.a tu taɾˈxeta ku ˈese ˈele ˈbi.a xe ˈuno ˈxota ku ˈatʃe} \\
8.12 & My WACRAL / GQRP / DXCC / FISTS number is 1234. & Mi número WACRAL / GQRP / DXCC / FISTS es 1224. & \ipa{mi ˈnumeɾo ˈdoble ˈube a θe ˈere ˈele / xe ku ˈere pe / de ˈekis θe θe / fɪsts es mil dosˈθjentos ˌbejntiˈkwatɾo} \\
\rowcolor{band} 9.0 & \textbf{Final Greetings} & \multicolumn{2}{l}{\textbf{Saludos Finales}} \\*
9.1 & Thank you. & Gracias. & \ipa{ˈɡɾaθjas} \\
9.2 & Thank you for a very enjoyable contact. & Gracias por tu contacto muy agradable. & \ipa{ˈɡɾaθjas poɾ tu konˈtakto muj aɡɾaˈdable} \\
9.3 & I enjoyed our chat. & Me gustó nuestra charla. & \ipa{me ɡusˈto ˈnwestɾa ˈtʃaɾla} \\
9.4 & I enjoyed our QSO in Spanish. & Me gustó mucho nuestro QSO en español. & \ipa{me ɡusˈto ˈmutʃo ˈnwestɾo ku ˈese o en espaˈɲol} \\
9.5 & It has been great fun. & Ha sido muy divertido. & \ipa{a ˈsido muj dibeɾˈtido} \\
9.6 & If you would like a regular contact with me, you can contact me via email. & Si quieres mantener un contacto regular conmigo, puedes contactarme vía email. & \ipa{si ˈkjeɾes manteˈneɾ un konˈtakto reɡuˈlaɾ konˈmiɡo ˈpwedes kontakˈtaɾme ˈbi.a ˈiːmeɪl} \\
9.7 & My email address is... & Mi dirección de email es… & \ipa{mi diɾekˈθjon de ˈiːmeɪl es …} \\
9.8 & My email address is shown on QSL.com / EHam.net & Mi dirección de email está en QSL.com / EHam.net. & \ipa{mi diɾekˈθjon de ˈiːmeɪl esˈta en ku ˈese ˈele ˈpunto kom / i hæm ˈpunto net} \\
9.9 & I look forward to meeting you again. & Espero encontrarte otra vez en el futuro. & \ipa{esˈpeɾo enkonˈtɾaɾte ˈotɾa beθ en el fuˈtuɾo} \\
9.11 & My best wishes to you and your family. & Mis mejores deseos para ti y para tu familia. & \ipa{mis meˈxoɾes deˈse.os ˈpaɾa ti i ˈpaɾa tu faˈmilja} \\
9.12 & Best wishes to your wife. & Los mejores deseos para tu esposa. & \ipa{los meˈxoɾes deˈse.os ˈpaɾa tu esˈposa} \\
9.121 & Best wishes to your husband. & Los mejores deseos para tu marido. & \ipa{los meˈxoɾes deˈse.os ˈpaɾa tu maˈɾido} \\
9.13 & Good luck. & Buena suerte. & \ipa{ˈbwena ˈsweɾte} \\
9.14 & Good health. & Buena salud. & \ipa{ˈbwena saˈlud} \\
9.15 & Goodbye & Adiós. & \ipa{aˈdjos} \\
9.16 & I wish you an enjoyable afternoon. & Te deseo una agradable tarde (de 12 a 18 horas). & \ipa{te deˈse.o ˈuna aɡɾaˈdable ˈtaɾde (de ˈdoθe a djeˈθotʃo ˈoɾas)} \\
9.161 & I wish you an enjoyable evening. & Te deseo una agradable tarde (después de las 18 horas). & \ipa{te deˈse.o ˈuna aɡɾaˈdable ˈtaɾde (desˈpwes de las djeˈθotʃo ˈoɾas)} \\
9.162 & I wish you an enjoyable weekend. & Te deseo un agradable fin de semana. & \ipa{te deˈse.o un aɡɾaˈdable fin de seˈmana} \\
9.17 & Enjoy your holiday. & Que te diviertas en tus vacaciones. & \ipa{ke te diˈbjeɾtas en tus bakaˈθjones} \\
9.18 & Happy New Year. & Feliz Año Nuevo. & \ipa{feˈliθ ˈaɲo ˈnwebo} \\
\rowcolor{band} 10.0 & \textbf{Country} & \multicolumn{2}{l}{\textbf{País}} \\*
10.1 & Albania & Albania & \ipa{alˈbanja} \\
10.2 & Austria & Austria & \ipa{ˈawstɾja} \\
10.3 & Belgium & Bélgica & \ipa{ˈbelxika} \\
10.4 & Belorussia & Bielorrusia & \ipa{bjeloˈrusja} \\
10.5 & Bosnia \& Herzegovina & Bosnia y Herzegovina & \ipa{ˈbosnja i eɾθeɡoˈbina} \\
10.6 & Bulgaria & Bulgaria & \ipa{bulˈɡaɾja} \\
10.7 & Croatia & Croacia & \ipa{kɾoˈaθja} \\
10.8 & Czech Republic & República Checa & \ipa{reˈpublika ˈtʃeka} \\
10.9 & Denmark & Dinamarca & \ipa{dinaˈmaɾka} \\
10.11 & Eire & Irlanda & \ipa{iɾˈlanda} \\
10.12 & England & Inglaterra & \ipa{inɡlaˈtera} \\
10.13 & Estonia & Estonia & \ipa{esˈtonja} \\
10.14 & Finland & Finlandia & \ipa{finˈlandja} \\
10.15 & France & Francia & \ipa{ˈfɾanθja} \\
10.16 & Germany & Alemania & \ipa{aleˈmanja} \\
10.17 & Gibraltar & Gibraltar & \ipa{xibɾalˈtaɾ} \\
10.18 & Greece & Grecia & \ipa{ˈɡɾeθja} \\
10.19 & Holland & Holanda & \ipa{oˈlanda} \\
10.21 & Hungary & Hungría & \ipa{uŋˈɡɾi.a} \\
10.22 & Iceland & Islandia & \ipa{isˈlandja} \\
10.23 & Northern Ireland & Irlanda del Norte & \ipa{iɾˈlanda del ˈnoɾte} \\
10.24 & Isle of Man & Isla de Man & \ipa{ˈisla de man} \\
10.25 & Italy & Italia & \ipa{iˈtalja} \\
10.26 & Latvia & Letonia & \ipa{leˈtonja} \\
10.27 & Liechtenstein & Liechtenstein & \ipa{ˈliçtn̩ʃtaɪn} \\
10.28 & Lithuania & Lituania & \ipa{liˈtwanja} \\
10.29 & Luxembourg & Luxemburgo & \ipa{luksemˈbuɾɡo} \\
10.31 & Macedonia & Macedonia & \ipa{maθeˈdonja} \\
10.32 & Moldavia & Moldavia & \ipa{molˈdabja} \\
10.33 & Netherlands & Países Bajos & \ipa{paˈises ˈbaxos} \\
10.34 & Norway & Noruega & \ipa{noˈɾweɡa} \\
10.35 & Poland & Polonia & \ipa{poˈlonja} \\
10.36 & Portugal & Portugal & \ipa{poɾtuˈɡal} \\
10.37 & Romania & Rumania & \ipa{ruˈmanja} \\
10.38 & Russian Federation & Federación Rusa & \ipa{fedeɾaˈθjon ˈrusa} \\
10.39 & Scotland & Escocia & \ipa{esˈkoθja} \\
10.41 & Slovakia & Eslovaquia & \ipa{esloˈbakja} \\
10.42 & Slovenia & Eslovenia & \ipa{esloˈbenja} \\
10.43 & Spain & España & \ipa{esˈpaɲa} \\
10.44 & Svalbard & Svalbard & \ipa{ˈsvalbard} \\
10.45 & Sweden & Suecia & \ipa{ˈsweθja} \\
10.46 & Switzerland & Suiza & \ipa{ˈswiθa} \\
10.47 & Turkey & Turquía & \ipa{tuɾˈki.a} \\
10.48 & Ukraine & Ucrania & \ipa{uˈkɾanja} \\
10.49 & Wales & Gales & \ipa{ˈɡales} \\
\end{longtable}
\normalsize
\section*{Hinweise zur Übertragung und Aussprache}
\fontsize{8.5}{10.5}\selectfont\raggedright
\textbf{Vorlage und Umfang.} Grundlage ist „Spanish v 1.1.doc“; Credits: Vicente -- EA5OO und
Colin Bass (ohne Rufzeichen). Alle 142 Phrasen und Ländernamen sind enthalten. Die ursprüngliche Gliederung
mit ihren Nummernsprüngen und dem fehlenden Abschnitt 7 bleibt erhalten.

\smallskip
\textbf{Dokumentierte Korrekturen.} „indicative“ wurde zu „indicativo“ (2.12), „antenna“ zu „antena“ (3.4--3.48),
„much calor“ zu „mucho calor“ (6.5), „favour“ zu „favor“ (8.6), „enviae“ zu „envía“ (8.11)
und „My dirección“ zu „Mi dirección“ (9.7). In 9.17 wurde „Que te diviertes“ zu „Que te diviertas“ berichtigt.
Fehlende Akzente wurden ergänzt: días, Qué, Cómo, frío, envío, envíe, Cuál, envíame, está, Adiós, Hungría
und Países. „Necesito“ nach dem Komma (5.2) und „Hay“ (6.12) wurden kleingeschrieben; mehrfache Leerzeichen
vereinheitlicht. Die fehlerhaften Webschreibweisen „QSLcon“ (8.8) und „QSLcom“ (9.8) wurden entsprechend
der englischen Spalte als „QSL.com“ wiedergegeben. Die IPA-Spalte folgt jeweils dem korrigierten spanischen Text.

\smallskip
\textbf{Abweichungen und Eigenheiten der Quelle.} In 8.12 steht englisch die Mitgliedsnummer 1234,
spanisch 1224. Beide Zahlen bleiben unverändert; die spanische Lautschrift liest „mil doscientos veinticuatro“.
„Rumania“ bleibt in 10.37 als gültige spanische Variante erhalten. Die Zeitangaben in Klammern bei 1.32,
1.33, 9.16 und 9.161 werden übernommen und ebenfalls transkribiert; sie können beim Funkgespräch entfallen.
Auch „Good evening“ wird in der Vorlage mit „Buenas tardes“ übersetzt; diese Zuordnung bleibt erhalten. Überwiegend steht die vertraute Anrede „tú“; „envíe“ in 8.3 entspricht der höflichen
Anrede „usted“. Für eine Sprecherin lauten „contento“ und „jubilado“: „contenta“ und „jubilada“.

\smallskip
\textbf{Transkriptionskonvention.} Die IPA-Spalte ist eine breite, manuell erstellte Aussprachehilfe
für eine gebräuchliche Aussprache Spaniens mit distinción und yeísmo, keine enge Transkription einer Aufnahme.
\ipa{θ} steht für z sowie c vor e/i; \ipa{s} für s. ll und konsonantisches y sind als \ipa{ʝ} wiedergegeben.
\ipa{ɾ} bezeichnet den einfachen Zungenschlag, \ipa{r} das gerollte r. b und v sind beide als \ipa{b}
transkribiert. Die kontextabhängigen weicheren Varianten von b, d und g (\ipa{β, ð, ɣ}) werden nicht eigens
markiert. \ipa{ˈ} zeigt die Hauptbetonung; \ipa{ˌ} kann in längeren Zahlen und Wörtern eine Nebenbetonung markieren.
Der Punkt markiert ausgewählte Vokalsilben. Satzmelodie und Wortverbindungen sind nicht dargestellt.

\smallskip
\textbf{Lateinamerikanische Aussprache.} In den meisten lateinamerikanischen Varianten sowie in Teilen Spaniens
werden z und c vor e/i wie s gesprochen (seseo): Ersetzen Sie dann in spanischen Wörtern \ipa{θ} durch \ipa{s}.
So wird „gracias“ zu \ipa{ˈɡɾasjas} und „cinco“ zu \ipa{ˈsinko}. Weitere regionale Unterschiede sind nicht dargestellt.

\smallskip
\textbf{Namen, Kürzel und Zahlen.} Peter, Sittingbourne und Dover sind mit englischer Aussprache wiedergegeben,
Hellschreiber und Liechtenstein mit deutscher. Icom ist mit einer englischen Aussprachevariante angegeben,
Loop und E-Mail ebenfalls englisch; „bureau“ französisch. Rufzeichen und die meisten Kürzel werden mit
spanischen Buchstabennamen gesprochen, nicht im internationalen Buchstabieralphabet: W „doble uve“,
Y hier „i griega“, V „uve“. FISTS wird als Wort, das e in eQSL und EHam als englisches \ipa{i} gelesen.
Zahlen: 706 „setecientos seis“, 599 „cinco nueve nueve“, 3.558 MHz
„tres coma cinco cinco ocho megahercios“, 8º C „ocho grados Celsius“.
Schrägstriche trennen Alternativen oder Listenelemente; Platzhalter \texttt{<\#\#\#\#\#>}
erscheinen in der Lautschrift als Auslassung und werden nicht mitgesprochen.

\smallskip
\textbf{LaTeX.} UTF-8-Datei; zweimal mit XeLaTeX kompilieren. Benötigt werden DejaVu Sans und die
üblichen LaTeX-Pakete der Präambel. IPA ist direkt als Unicode eingegeben; keine externen Dateien erforderlich.
\end{document}
